\section{\sst{}}
\label{sec:results-sst2}

\section{\sst{}}
\label{sec:results-sst2}

\section{\sst{}}
\label{sec:results-sst2}

\section{\sst{}}
\label{sec:results-sst2}

\input{tables/nlp_sst2}

\Cref{tab:sst2} reports the main \sst{} fine-tuning results. Flat \adamw{} reaches $89.56\%$ peak validation accuracy and $88.00\%$ final accuracy. Adding warmup-linear scheduling improves this to $90.48\%$ peak accuracy and $89.84\%$ final accuracy, confirming that the scheduler is a strong baseline in this setting.

\method{} improves over flat \adamw{}, reaching $91.03\%$ peak accuracy and $90.25\%$ final accuracy. When combined with warmup-linear scheduling, \method{} reaches a lower peak accuracy of $90.80\%$, but a higher final accuracy of $90.44\%$ and a smaller best-to-final drop of $0.37$ percentage points. This suggests that the scheduled variant is more stable, even though the unscheduled variant reaches a slightly higher validation peak.

The axis ablation in \Cref{tab:sst2-axis} clarifies which adaptive component is most useful. \method{}-Noise with warmup-linear scheduling gives the strongest result, reaching $91.12\%$ peak accuracy and $90.57\%$ final accuracy. In comparison, \method{}-LR with warmup-linear scheduling reaches $90.55\%$ peak and $90.21\%$ final accuracy, while \method{}-Dual reaches $90.80\%$ peak and $90.44\%$ final accuracy. Thus, adaptive gradient-noise control is the strongest individual axis in this setting, and combining it with adaptive learning-rate modulation does not improve over the noise-only variant.

Overall, the \sst{} results provide moderate evidence that \method{} can improve Transformer fine-tuning, but the gain is smaller than in noisy \agnews{}. The main conclusion is not that all adaptive axes are useful, but that adaptive gradient-noise control provides the most consistent benefit among the tested \method{} variants.

\begin{figure}[t]
    \centering
    \thesisfigure{0.90\textwidth}{figures/nlp/sst2_curves.pdf}
    \caption[
        Validation trajectories on \sst{}
    ]{
        Validation trajectories on \sst{}. The \method{} variants improve over the standard warmup-linear baseline in peak or final validation accuracy, with the noise-only variant giving the strongest peak and final accuracy among the tested axis-ablation configurations.
    }
    \label{fig:sst2-curves}
\end{figure}

\Cref{tab:sst2} reports the main \sst{} fine-tuning results. Flat \adamw{} reaches $89.56\%$ peak validation accuracy and $88.00\%$ final accuracy. Adding warmup-linear scheduling improves this to $90.48\%$ peak accuracy and $89.84\%$ final accuracy, confirming that the scheduler is a strong baseline in this setting.

\method{} improves over flat \adamw{}, reaching $91.03\%$ peak accuracy and $90.25\%$ final accuracy. When combined with warmup-linear scheduling, \method{} reaches a lower peak accuracy of $90.80\%$, but a higher final accuracy of $90.44\%$ and a smaller best-to-final drop of $0.37$ percentage points. This suggests that the scheduled variant is more stable, even though the unscheduled variant reaches a slightly higher validation peak.

The axis ablation in \Cref{tab:sst2-axis} clarifies which adaptive component is most useful. \method{}-Noise with warmup-linear scheduling gives the strongest result, reaching $91.12\%$ peak accuracy and $90.57\%$ final accuracy. In comparison, \method{}-LR with warmup-linear scheduling reaches $90.55\%$ peak and $90.21\%$ final accuracy, while \method{}-Dual reaches $90.80\%$ peak and $90.44\%$ final accuracy. Thus, adaptive gradient-noise control is the strongest individual axis in this setting, and combining it with adaptive learning-rate modulation does not improve over the noise-only variant.

Overall, the \sst{} results provide moderate evidence that \method{} can improve Transformer fine-tuning, but the gain is smaller than in noisy \agnews{}. The main conclusion is not that all adaptive axes are useful, but that adaptive gradient-noise control provides the most consistent benefit among the tested \method{} variants.

\begin{figure}[t]
    \centering
    \thesisfigure{0.90\textwidth}{figures/nlp/sst2_curves.pdf}
    \caption[
        Validation trajectories on \sst{}
    ]{
        Validation trajectories on \sst{}. The \method{} variants improve over the standard warmup-linear baseline in peak or final validation accuracy, with the noise-only variant giving the strongest peak and final accuracy among the tested axis-ablation configurations.
    }
    \label{fig:sst2-curves}
\end{figure}

\Cref{tab:sst2} reports the main \sst{} fine-tuning results. Flat \adamw{} reaches $89.56\%$ peak validation accuracy and $88.00\%$ final accuracy. Adding warmup-linear scheduling improves this to $90.48\%$ peak accuracy and $89.84\%$ final accuracy, confirming that the scheduler is a strong baseline in this setting.

\method{} improves over flat \adamw{}, reaching $91.03\%$ peak accuracy and $90.25\%$ final accuracy. When combined with warmup-linear scheduling, \method{} reaches a lower peak accuracy of $90.80\%$, but a higher final accuracy of $90.44\%$ and a smaller best-to-final drop of $0.37$ percentage points. This suggests that the scheduled variant is more stable, even though the unscheduled variant reaches a slightly higher validation peak.

The axis ablation in \Cref{tab:sst2-axis} clarifies which adaptive component is most useful. \method{}-Noise with warmup-linear scheduling gives the strongest result, reaching $91.12\%$ peak accuracy and $90.57\%$ final accuracy. In comparison, \method{}-LR with warmup-linear scheduling reaches $90.55\%$ peak and $90.21\%$ final accuracy, while \method{}-Dual reaches $90.80\%$ peak and $90.44\%$ final accuracy. Thus, adaptive gradient-noise control is the strongest individual axis in this setting, and combining it with adaptive learning-rate modulation does not improve over the noise-only variant.

Overall, the \sst{} results provide moderate evidence that \method{} can improve Transformer fine-tuning, but the gain is smaller than in noisy \agnews{}. The main conclusion is not that all adaptive axes are useful, but that adaptive gradient-noise control provides the most consistent benefit among the tested \method{} variants.

\begin{figure}[t]
    \centering
    \thesisfigure{0.90\textwidth}{figures/nlp/sst2_curves.pdf}
    \caption[
        Validation trajectories on \sst{}
    ]{
        Validation trajectories on \sst{}. The \method{} variants improve over the standard warmup-linear baseline in peak or final validation accuracy, with the noise-only variant giving the strongest peak and final accuracy among the tested axis-ablation configurations.
    }
    \label{fig:sst2-curves}
\end{figure}

\Cref{tab:sst2} reports the main \sst{} fine-tuning results. Flat \adamw{} reaches $89.56\%$ peak validation accuracy and $88.00\%$ final accuracy. Adding warmup-linear scheduling improves this to $90.48\%$ peak accuracy and $89.84\%$ final accuracy, confirming that the scheduler is a strong baseline in this setting.

\method{} improves over flat \adamw{}, reaching $91.03\%$ peak accuracy and $90.25\%$ final accuracy. When combined with warmup-linear scheduling, \method{} reaches a lower peak accuracy of $90.80\%$, but a higher final accuracy of $90.44\%$ and a smaller best-to-final drop of $0.37$ percentage points. This suggests that the scheduled variant is more stable, even though the unscheduled variant reaches a slightly higher validation peak.

The axis ablation in \Cref{tab:sst2-axis} clarifies which adaptive component is most useful. \method{}-Noise with warmup-linear scheduling gives the strongest result, reaching $91.12\%$ peak accuracy and $90.57\%$ final accuracy. In comparison, \method{}-LR with warmup-linear scheduling reaches $90.55\%$ peak and $90.21\%$ final accuracy, while \method{}-Dual reaches $90.80\%$ peak and $90.44\%$ final accuracy. Thus, adaptive gradient-noise control is the strongest individual axis in this setting, and combining it with adaptive learning-rate modulation does not improve over the noise-only variant.

Overall, the \sst{} results provide moderate evidence that \method{} can improve Transformer fine-tuning, but the gain is smaller than in noisy \agnews{}. The main conclusion is not that all adaptive axes are useful, but that adaptive gradient-noise control provides the most consistent benefit among the tested \method{} variants.

\begin{figure}[t]
    \centering
    \thesisfigure{0.90\textwidth}{figures/nlp/sst2_curves.pdf}
    \caption[
        Validation trajectories on \sst{}
    ]{
        Validation trajectories on \sst{}. The \method{} variants improve over the standard warmup-linear baseline in peak or final validation accuracy, with the noise-only variant giving the strongest peak and final accuracy among the tested axis-ablation configurations.
    }
    \label{fig:sst2-curves}
\end{figure}